% Research note only. Exact incidence calculation; conditional object budget.
% Not a numbered manuscript theorem, not an unbounded-period count, no movement.
\subsection*{The precise generation hypothesis for an intrinsic order atlas}

This note isolates the remaining input after the rank-period generating-anchor
construction and the qualitative class-openness theorem.  Let $D$ be a proper
one-dimensional scheme over an algebraically closed field $k$, possibly
nonreduced.  Let $\mathcal B$ be a finite locally free unital
$\mathcal O_D$-algebra of positive rank, and let $E$ be locally isomorphic to
$\mathcal B^s$ as a right $\mathcal B$-module.  In the proposed geometric
application, $\operatorname{rk}_{\mathcal O_D}\mathcal B=M^2$,
$s=r/M$, and $D$ is a fixed multiple of the inherited Cartier cycle.

\paragraph{An exact incidence statement.}
Assume that the underlying $\mathcal O_D$-module $E$ is globally generated.
Choose a finite-dimensional space $W\subset H^0(D,E)$ whose evaluation is
surjective.  Then $s+1$ general elements of $W$ generate $E$ as a
$\mathcal B$-module.

Here is the complete calculation.  At a closed point $x$, write
\[
 \mathcal B(x)/\operatorname{rad}\mathcal B(x)
       =\prod_i\operatorname{Mat}_{d_i}(k).
\]
For $h$ candidate generators, surjectivity is equivalent, by Nakayama over
the finite-dimensional algebra, to surjectivity on each of these blocks.
On the $i$th block, the map from $\mathcal B(x)^h$ to
$\mathcal B(x)^s$ is an arbitrary $(s d_i)\times(h d_i)$ matrix.
The locus where it is not surjective has codimension
\[
                         d_i(h-s)+1.
\]
Evaluation $W^h\to E(x)^h$ is surjective as a map of vector spaces, so this
is also the codimension of the corresponding inverse image in $W^h$.
For $h=s+1$ it is at least two.  The universal failure locus is closed in
$D\times W^h$, because it is the support of the cokernel of the actual
universal evaluation map.  Its dimension is at most
$\dim W^h-2+\dim D=\dim W^h-1$.  Properness of $D$ makes its image in
$W^h$ closed; the dimension bound makes that image proper.  Its complement
therefore parametrizes actual surjections
\[
                 \mathcal B^{s+1}\twoheadrightarrow E.
\]
The argument uses the closed fibres of the nonreduced curve; Nakayama then
gives surjectivity on the whole curve, including its nilpotent structure.

The kernel is locally free of rank one over $\mathcal B$.  Indeed, over a
local ring of the centre the algebra is semilocal.  The displayed surjection
onto the free module $\mathcal B^s$ splits, so its kernel is projective and
stably free.  Its reduction modulo the Jacobson radical is the regular
module over $\mathcal B/\operatorname{rad}\mathcal B$.  Lifting one such
generator and applying Nakayama gives a surjection $\mathcal B\to K$.
It splits because $K$ is projective; the remaining summand has zero
semisimple reduction and is zero by Nakayama.  Thus $K\simeq\mathcal B$
locally.  The same cancellation argument applies to an actual surjection
between such locally free order-module families over any coefficient ring;
it is not a claim that arbitrary saturated module families are flat.

\paragraph{Why order generation alone is insufficient.}
The surjectivity of $W\to E(x)$ as a vector-space map is essential.  Already
for $B=\operatorname{Mat}_d(k)$ and its regular right module $E=B$, the
$d$ matrices $e_i e_1^{\mathsf T}$ generate $E$ as a right $B$-module.
But $h$ scalar linear combinations of these matrices have all their column
images in a common subspace of $k^d$ of dimension at most $h$.
Their right-$B$ span cannot equal $B$ if $h<d$.
In particular two scalar combinations need not suffice even though $s=1$.
Replacing ``underlying ordinary module is globally generated'' by
``a given list generates over $B$'' invalidates the incidence proof.

The stronger geometric obstruction
\ref{D:elliptic-cone-order-generators} is also still relevant: for arbitrary
anchors on one fixed cone, the actual module-generator count is $3h$ even
at generic order rank one and with Frobenius-periodic full modules.  Its
Brauer class is zero, so it does not rule out a better, suitably normalized
rank-period anchor.  It does rule out deriving the required bound from
generic Morita rank alone.

\paragraph{A conditional small-presentation budget.}
Suppose the following genuinely geometric statements have been established
for the chosen rank-$M$ anchors and every geometric unit object being
counted.  They are hypotheses here, not consequences of the incidence
calculation or of a conductor sandwich.

\begin{enumerate}
\item After a fixed twist that preserves a fixed bound on the actual unit
      poles, the underlying ordinary $E$ is globally generated.  The
      resulting rank-one order kernel admits two-generator presentations
      after twists of size $O(s)$ by the fixed Cartier cycle.  Iterating a
      fixed number of times gives the separated partial resolutions used
      for exact low-degree Hom and composition.
\item The reference terms in those partial resolutions are fixed twists of
      $\mathcal B$: the degree-zero term has multiplicity $s+1$, while each
      negative-degree term has multiplicity two.  Kernels may vary, but the
      reference terms do not.  Actual local split exactness on the retained
      finite cycle is imposed, so the quotient and last syzygy are locally
      free order-modules over every coefficient ring.
\item All fixed coefficient complexes required for these terms, their
      Frobenius correspondences and opposite correspondences have the
      following bounds: the untwisted top--top block has total cohomology
      dimension $O_R(tM^2)$, and every block involving a lower term has
      total dimension $O_R(tM^2(1+s))$.  This is a bound on the actual fixed
      complexes, not just on generic ranks or virtual Euler characteristics.
\item The retained cycle and exact Morita-unit comparison have fixed
      Cartier losses independent of $M,s,t$.  An $O(s)$ twist in a reference
      syzygy is not substituted for the pole bound of the actual unit object.
\end{enumerate}

Under these hypotheses the repeated-term accounting is explicit.  Intrinsic
Hom between two free order terms is a twist of $\mathcal B$, rather than an
arbitrary ordinary Hom of rank $M^4$.  The top--top multiplicity is
$(s+1)^2$, a mixed multiplicity is $O(s+1)$, and a bottom--bottom
multiplicity is $O(1)$.  There are only a fixed number of blocks.  Hence
their total finite-complex size is
\[
 O_R\bigl(tM^2(s+1)^2
       +tM^2(1+s)(s+1)+tM^2(1+s)\bigr)
                         =O_R(tr^2),\qquad r=Ms,\ s\ge1.
\]
Fixed $k$-linear contractions of these coefficient complexes commute with
arbitrary coefficient extension.  The previously proved bounded-width
Hom perturbation construction then has degree bounded independently of
$M,s,t$ (with the usual dependence on the one-step $p$-Frobenius).
It supplies actual cycles, boundaries, compositions and inverse-product
witnesses with $O_R(ftr^2)$ coordinates over $f$ phases, once the stated
exact order-module presentation and finite-effectivity hypotheses hold.
The separated source/target lengths and the top-generator identity test
must be retained; equal-depth truncated resolutions can introduce spurious
bottom-degree maps.

This conditional accounting introduces neither arbitrary matrices for the
action of every algebra generator nor bases for varying cohomology groups.
It therefore identifies a sufficient geometric presentation input, but
does not establish that input.

\paragraph{What the newly proved comparison theorem removes.}
Theorem-level statement
\ref{D:effective-formal-unit-family-class-openness} applies to an already
effective finite-type unit-family.  Its qualitative Serre quotient and
exceptional Picard rigidity allow comparison bounds depending on the chart
and the closed class, without a quantitative charge to object coordinates.
Thus a separate fixed-determinant enumeration is not required merely to
prove class openness for a successful intrinsic atlas.  The theorem does
not supply the geometric generation and syzygy estimates above, and does
not change the original object Frobenius poles.

The exact current obstruction is therefore the first three geometric
presentation hypotheses, with their compatibility with the fourth.  The
canonical comparisons $P\subset B\subset P(C)$ and
$Q\subset E\subset Q(C)$ do not by themselves provide these assertions.
No unbounded-geometric-period hereditary count, group-generation result,
or new main-gate movement is claimed in this note.
